% Candidate source-aware packet for production Rethlas.
% Do not insert into a live chapter before verification.
% Do not transmit without the user's explicit authorization.
% No novelty or priority claim is intended.

\paragraph{Imported inputs.}
We use the following results with their stated hypotheses.
\begin{enumerate}
\item Chambert--Loir--Ducros, Definitions~8.1.1 and~8.1.13 and
Lemma~8.1.4: the CLD psh sheaf is the smallest subsheaf containing the
smooth functions with nonnegative $\mathrm d'\mathrm d''$ and stable under
decreasing sequential limits; its first stage consists of the locally
decreasing limits of smooth psh functions.
\item Burgos Gil--Gubler--Jell--K\"unnemann, Definition~3.1.2,
Proposition~3.1.3, Theorem~3.1.4, and Theorem~3.3.4: on the tropical toric
partial compactification $N_\Sigma$, their psh condition is local and stable
under decreasing sequences, is equivalent to continuity plus convexity on
the dense stratum, and is equivalent to being locally a decreasing limit of
smooth psh functions.  Their smooth structure is the standard Lagerberg
smooth structure on $N_\Sigma$.
\item The elementary uniform-to-monotone device of Chambert--Loir--Ducros,
Lemma~8.1.10: if real functions $f_j$ converge uniformly to $f$, then
$f_j+2\lVert f_j-f\rVert_\infty$ has a decreasing subsequence converging to
$f$.
\item Farkas' lemma in the following finite-dimensional form.  For vectors
$v_1,\ldots,v_m\in\mathbf R^N$ and $A=(A_i)\in\mathbf R^m$, there is a
covector $\xi\in(\mathbf R^N)^*$ with $\xi(v_i)\leq A_i$ for every $i$ if
and only if
\[
 \sum_i\lambda_iA_i\geq0
 \quad\text{for every }\lambda_i\geq0
 \text{ such that }\sum_i\lambda_i v_i=0.
\]
\item A smooth regularized maximum $M_\varepsilon$ of finitely many real
variables can be chosen convex and nondecreasing in every variable, with
\[
 \max_j s_j\leq M_\varepsilon(s)\leq\max_j s_j+\varepsilon,
\]
and with the truncation property that it equals one entry locally whenever
that entry exceeds all the others by at least $2\varepsilon$.
\end{enumerate}

\begin{proposition}
\label{prop:p07-toric-cld-collapse-candidate}
Let $N$ be a lattice, let $\Sigma$ be a rational polyhedral fan in
$N_{\mathbf R}$, and equip the tropical toric partial compactification
$N_\Sigma$ with its standard Lagerberg smooth structure.  On every open
subset of $N_\Sigma$, the following three sheaves coincide:
\[
 \operatorname{Psh}^{\mathrm{CLD}}_{N_\Sigma,1},\qquad
 \operatorname{Psh}^{\mathrm{CLD}}_{N_\Sigma},\qquad
 \operatorname{Psh}^{\mathrm{tor}}_{N_\Sigma},
\]
where the last sheaf is the psh sheaf of Burgos
Gil--Gubler--Jell--K\"unnemann.  In particular, the CLD hierarchy on
$N_\Sigma$ collapses at its first decreasing-limit stage.
\end{proposition}

\begin{proof}
Let $f$ be a standard Lagerberg-smooth function with
$\mathrm d'\mathrm d''f\geq0$.  It is continuous, and its restriction to
the dense stratum $N_{\mathbf R}$ is convex because its Hessian is positive
semidefinite.  Burgos Gil--Gubler--Jell--K\"unnemann, Theorem~3.1.4,
therefore puts $f$ in $\operatorname{Psh}^{\mathrm{tor}}$.  Their
Proposition~3.1.3 says that $\operatorname{Psh}^{\mathrm{tor}}$ is a sheaf
and is stable under decreasing sequential limits.  The defining minimality
of the CLD sheaf gives
\[
 \operatorname{Psh}^{\mathrm{CLD}}_{N_\Sigma}
 \subseteq \operatorname{Psh}^{\mathrm{tor}}_{N_\Sigma}.
\]

Conversely, their Theorem~3.3.4 writes every section of
$\operatorname{Psh}^{\mathrm{tor}}$ locally as a decreasing sequence of
smooth psh functions for the standard Lagerberg smooth structure.  These
are precisely the generators used at CLD stage zero, so
\[
 \operatorname{Psh}^{\mathrm{tor}}_{N_\Sigma}
 \subseteq \operatorname{Psh}^{\mathrm{CLD}}_{N_\Sigma,1}.
\]
The inclusion of the first CLD stage in the generated CLD sheaf is
tautological.  The three sheaves are equal.
\end{proof}

\begin{theorem}
\label{thm:p07-finite-tropical-curve-collapse-candidate}
Let $X$ be a one-dimensional boundaryless affine tropical space.  Assume
that every point has a neighbourhood isomorphic, as a calibrated affine
tropical space with smooth structure, to an open subset of the support of a
finite balanced one-dimensional polyhedral complex in some $\mathbf R^N$,
with the smooth structure induced by ambient smooth functions.  Then
\[
                    \operatorname{Psh}_{X,1}
                    =\operatorname{Psh}_X.
\]
Equivalently, every CLD-psh function on $X$ is locally a decreasing limit of
smooth CLD-psh functions.
\end{theorem}

\begin{proof}
The assertion is local.  At a point in the interior of an edge it is the
usual one-dimensional convex regularization.  It remains to work at a
vertex.  After shrinking, write the local star as
\[
 S=\{0\}\cup\coprod_{i=1}^m(0,r_i)v_i\subset\mathbf R^N,
 \qquad
 \sum_{i=1}^m w_i v_i=0,
 \qquad w_i>0.
\]
For a function $u:S\to\mathbf R\cup\{-\infty\}$, put
$u_i(t)=u(tv_i)$.

We first record two conditions satisfied by every CLD-psh function on $S$:
each $u_i$ is convex in the extended-real sense, and
\begin{equation}
\label{eq:p07-junction-jensen}
 u(0)\leq\sum_{i=1}^m\alpha_i u_i(t_i)
\end{equation}
whenever $\alpha_i\geq0$, $\sum_i\alpha_i=1$,
$0\leq t_i<r_i$, and $\sum_i\alpha_i t_i v_i=0$.
Indeed, let first $u$ be smooth psh.  Its restriction to every branch is
convex.  If
\[
                   \ell(x)=u(0)+\mathrm d u_0(x)
\]
is its ambient affine tangent at the vertex, branchwise convexity gives
$u_i(t)\geq\ell(tv_i)$.  Taking the indicated convex combination proves
\eqref{eq:p07-junction-jensen}.  The class of functions with these two
properties is local, contains the smooth psh functions, and is stable under
decreasing sequential limits: all inequalities pass to the limit because
the sums are finite and the functions never take the value $+\infty$.
It is also stable under the local unions used at limit ordinals.  The
minimal definition of $\operatorname{Psh}_X$ proves the claim for every
CLD-psh $u$.

We next dispose of the value $-\infty$.  Suppose first that $u(0)$ is
finite and set
\[
                  L_i=\lim_{t\downarrow0}u_i(t).
\]
Upper semicontinuity gives $L_i\leq u(0)$.  Apply
\eqref{eq:p07-junction-jensen} with a common small $t$ and with
\[
                 \alpha_i=\frac{w_i}{\sum_jw_j}.
\]
The balancing relation makes this an admissible convex combination.  On
letting $t\downarrow0$, we obtain
\[
 u(0)\leq\sum_i\alpha_iL_i\leq u(0).
\]
All $\alpha_i$ are positive, so $L_i=u(0)$ for every $i$.  If some
$u_i(t_0)$ were $-\infty$, convexity between $0$ and $t_0$ would force
$u_i(t)=-\infty$ for $0<t<t_0$, contradicting $L_i=u(0)$.  Thus $u$ is
finite and continuous on a smaller star.

If $u(0)=-\infty$, then on each branch either $u_i$ is already
identically $-\infty$ near zero, or there is a $t_i>0$ with $u_i(t_i)$
finite.  In the second case, convexity between $0$ and $t_i$ again gives
$u_i(t)=-\infty$ for $0<t<t_i$.  There are only finitely many branches, so
after one common shrinking $u\equiv-\infty$.  The smooth psh constants
$-k$ decrease to this germ.  Hence it remains only to treat the finite
continuous case.

Choose radii $0<\rho_i<R_i<r_i$ and put
\[
 V=\{0\}\cup\coprod_i(0,\rho_i)v_i,
 \qquad
 K=\{0\}\cup\coprod_i(0,R_i]v_i.
\]
Thus $V$ is one fixed open neighbourhood of the vertex and its closure is
contained in the compact star $K\Subset S$.  On each interval $[0,R_i]$
choose a finite partition starting at zero, and let $h_i$ be the chordal
interpolant of $u_i$.  Choose the partitions with mesh tending to zero.
Convexity and uniform continuity show that the resulting continuous
functions $h$ are piecewise affine and convex on each branch, satisfy
$h\geq u$, and converge uniformly to $u$ on $K$.

Let $t_i>0$ be the first positive partition point on branch $i$ and set
\[
                   A_i=\frac{u_i(t_i)-u(0)}{t_i}.
\]
We claim that there is a covector $\xi\in(\mathbf R^N)^*$ such that
\begin{equation}
\label{eq:p07-support-covector}
                         \xi(v_i)\leq A_i
                         \quad(1\leq i\leq m).
\end{equation}
By Farkas' lemma, it is enough to prove
$\sum_i\lambda_iA_i\geq0$ whenever $\lambda_i\geq0$ and
$\sum_i\lambda_i v_i=0$.  If $\lambda$ is nonzero, put
\[
 \beta_i=
 \frac{\lambda_i/t_i}{\sum_j\lambda_j/t_j}.
\]
Then $\sum_i\beta_i=1$ and $\sum_i\beta_it_iv_i=0$.  Applying
\eqref{eq:p07-junction-jensen} to the points $t_iv_i$, subtracting $u(0)$,
and clearing the positive denominator gives
\[
                  0\leq\sum_i\lambda_i
                  \frac{u_i(t_i)-u(0)}{t_i}
                  =\sum_i\lambda_iA_i.
\]
This proves \eqref{eq:p07-support-covector}.

Since the slopes of a convex chordal interpolant are nondecreasing, the
ambient affine function
\[
                         \ell(x)=u(0)+\xi(x)
\]
satisfies $\ell\leq h$ throughout $K$.  We now approximate $h$ by a smooth
psh function on the fixed open star $V$.  On branch $i$, write the convex
piecewise-affine function $h_i$ as the maximum of the affine extensions of
its finitely many pieces:
\[
                         h_i(t)=\max_{1\leq j\leq q_i}L_{ij}(t).
\]
For $C>0$, set
\[
                         B_i(t)=u(0)+C+\xi(tv_i).
\]
Choose $0<2\varepsilon<C$ and define
\[
 g_i(t)=M_\varepsilon
       \bigl(B_i(t),L_{i1}(t),\ldots,L_{iq_i}(t)\bigr).
\]
At $t=0$, the entry $B_i$ has value $u(0)+C$, while every $L_{ij}(0)$ is
at most $u(0)$.  The truncation property of the regularized maximum implies
that, on an initial interval of each branch, $g_i=B_i$.  Since there are
finitely many branches, the $g_i$ therefore glue near the vertex to the
single ambient affine function
\[
                         x\longmapsto u(0)+C+\xi(x).
\]
Away from the vertex they are ordinary smooth functions on open intervals.
Thus the glued function $g$ is smooth for the given tropical smooth
structure.  Each $g_i$ is convex, because the regularized maximum is convex
and nondecreasing and all its entries are affine.  Hence
$\mathrm d'\mathrm d''g\geq0$.

Finally, $\ell\leq h$ implies $B_i\leq h_i+C$, and the regularized-maximum
estimate gives
\[
                         h\leq g\leq h+C+\varepsilon
                         \quad\text{on }K.
\]
Choosing $C$ and $\varepsilon$ tending to zero proves that every chordal
approximant $h$ is uniformly approximable from above by functions smooth
and psh on the fixed open star $V$.  Indeed, the affine extensions in the
definition of $g_i$ define $g_i$ on all of $(0,\rho_i)$, and for each $g$
the truncation argument above gives ambient-affine agreement on some
neighbourhood of the vertex; the size of that latter neighbourhood may
depend on $g$, as is allowed by smoothness.

For the successive chordal interpolants $h_j$, choose smooth psh $g_j$ with
\[
                         \lVert g_j-h_j\rVert_K\leq j^{-1}.
\]
Then $g_j\to u$ uniformly on $K$.  Put
$\delta_j=\lVert g_j-u\rVert_K$.  By the uniform-to-monotone lemma, the
sequence $g_j+2\delta_j$ has a decreasing subsequence converging uniformly,
and hence pointwise, to $u$ on $K$.  In particular the same sequence is
defined, decreasing, and convergent to $u$ on the one fixed open
neighbourhood $V$.  Adding a constant preserves smooth
plurisubharmonicity.  This is exactly first-stage membership, so $u$ belongs
to $\operatorname{Psh}_{X,1}$ near the vertex.  The proof is complete.
\end{proof}

\paragraph{Boundary of the result.}
The curve theorem uses finiteness of the local star and a positive balancing
relation through all its branches.  It does not cover tropical boundary
points or the infinite-star spaces of Chambert--Loir--Ducros, Section~6.7.3.
In dimension at least two, the analogous Jensen reduction leaves a genuine
polyhedral regularization problem: one needs face-compatible
piecewise-affine convex approximation and then CLD-smooth approximation.
The regularity property of Botero--Mazzon--Pille-Schneider,
arXiv:2603.08664, isolates the first issue but is not proved for arbitrary
balanced polyhedral spaces.  No claim about the unrestricted hierarchy is
made here.
